\section*{Chapter \ref{ch:b3:innate-immunity} --- Innate Immunity and Inflammation}

\begin{solution}{exo:b3:innate-immunity:1}
\omterm{def:b3:innate-immunity:innate}{Barriers}: skin keratin (mechanical), airway mucus and cilia (trapping
and clearance), stomach acid (kills swallowed microbes), lysozyme and
defensins in secretions (lyse and puncture), the commensal \omterm{def:b3:microbiomes:symbiosis}{microbiome}
(occupies niches). Cells: \omterm{def:b3:innate-immunity:innate}{neutrophils} (kill bacteria by \omterm{def:b3:innate-immunity:killing}{phagocytosis}),
\omterm{def:b3:innate-immunity:innate}{macrophages} (eat, clean up, raise the alarm), \omterm{def:b3:innate-immunity:innate}{dendritic cells} (carry
antigen to lymph nodes and start adaptive responses), mast cells
(release histamine, inflame), \omterm{def:b3:innate-immunity:innate}{natural killer cells} (kill infected and
transformed cells).
\end{solution}

\begin{solution}{exo:b3:innate-immunity:2}
Shared by a whole class of microbes, essential to them (so it cannot
be lost to escape), and absent from the host. \omterm{def:b3:bacteriology:envelope}{Lipopolysaccharide} ---
TLR4; flagellin --- TLR5; double-stranded RNA --- TLR3 (and RIG-I in the
cytosol); unmethylated \omterm{def:b3:chromatin-epigenetics:methylation}{CpG} DNA --- TLR9; \omterm{def:b3:bacteriology:envelope}{peptidoglycan} fragments --- NOD
proteins; $\beta$-glucan --- dectin-1.
\end{solution}

\begin{solution}{exo:b3:innate-immunity:3}
Redness: arteriolar dilation by histamine, nitric oxide and
prostaglandins. Heat: the same increased blood flow. Swelling: venular
leak of plasma after endothelial contraction (histamine, C3a, C5a).
Pain: bradykinin and prostaglandin E$_{2}$ sensitising nociceptors,
plus pressure from the swelling.
\end{solution}

\begin{solution}{exo:b3:innate-immunity:4}
Classical: C1q binding antibody (or \omterm{prop:b3:innate-immunity:systemic}{C-reactive protein}) on a surface.
Lectin: mannose-binding lectin on microbial sugars. Alternative:
spontaneous C3 tick-over and C3b deposition on any unprotected surface.
Outcomes: C3b opsonises for \omterm{def:b3:innate-immunity:killing}{phagocytosis}; C3a and C5a inflame and
recruit; C5b starts the \omterm{def:b3:innate-immunity:complement}{membrane attack complex} that lyses.
\end{solution}

\begin{solution}{exo:b3:innate-immunity:5}
$B = [k_{0}/(a-d)](e^{(a-d)t} - 1)$ with $a - d = 0.1$: at \qty{60}{s},
$20(e^{6} - 1) \approx 8000$; at \qty{90}{s}, $20(e^{9} - 1) \approx
1.6\times 10^{5}$. Host cell: $k_{0}/(d - a) = 2/0.1 = 20$.
\end{solution}

\begin{solution}{exo:b3:innate-immunity:6}
Rolling on \omterm{def:b3:innate-immunity:inflammation}{selectins} (lectins binding carbohydrate); \omterm{def:b3:innate-immunity:inflammation}{chemokine}
signalling that activates \omterm{def:b3:innate-immunity:inflammation}{integrins}; firm adhesion of \omterm{def:b3:innate-immunity:inflammation}{integrins} to
endothelial adhesion molecules (ICAM); transmigration between
endothelial cells; chemotaxis along the gradient. Without
$\beta_{2}$ \omterm{def:b3:innate-immunity:inflammation}{integrins} the \omterm{def:b3:innate-immunity:innate}{neutrophils} roll but never stop or cross:
they pile up in the blood (very high counts), infections recur without
pus forming, wounds heal badly, and the umbilical cord separates late.
\end{solution}

\begin{solution}{exo:b3:innate-immunity:7}
Chronic granulomatous disease: the \omterm{def:b3:innate-immunity:killing}{NADPH oxidase} is defective, so no
superoxide, hydrogen peroxide or hypochlorite is made in the phagosome.
Infections by organisms that destroy their own hydrogen peroxide
(catalase-positive): \emph{Staphylococcus aureus}, \emph{Aspergillus},
\emph{Burkholderia}, \emph{Serratia}, \emph{Nocardia}. Granulomas form
because \omterm{def:b3:innate-immunity:innate}{macrophages} engulf microbes they cannot kill and wall them
off, with continual recruitment, in a chronic lesion.
\end{solution}

\begin{solution}{exo:b3:innate-immunity:8}
Inhibitory receptors count MHC~I; activating receptors count stress
ligands; the cell dies if activation outweighs inhibition. (a) Losing
MHC~I removes the inhibition: the \omterm{def:b3:innate-immunity:innate}{natural killer cell} kills it --- the
cost of escaping T~cells. (b) With MHC~I kept but stress ligands
displayed, activation can still win and the cell is killed; \omterm{def:b3:cancer-biology:hallmarks}{tumours}
therefore often shed or suppress the stress ligands.
\end{solution}

\begin{solution}{exo:b3:innate-immunity:9}
A pure protein carries no pathogen pattern, so the \omterm{def:b3:innate-immunity:innate}{dendritic cells} that
take it up do not mature, present it without costimulation, and the
T~cells that recognise it are switched off or ignore it. An adjuvant
--- alum, a lipid, a \omterm{def:b3:innate-immunity:prr}{Toll-like receptor} ligand --- triggers pattern
receptors, matures the \omterm{def:b3:innate-immunity:innate}{dendritic cell} and supplies the second signal.
This is Janeway's point: the adaptive system responds to antigen only
when the innate system says the antigen came with a microbe.
\end{solution}

\begin{solution}{exo:b3:innate-immunity:10}
Without fever, $48$ doublings: $10^{4}\times 2^{48} \approx 3\times
10^{18}$ (an absurdity that shows killing, not growth, decides the
outcome). With fever, $32$ doublings: $4\times 10^{13}$ --- a factor
$2^{16} \approx 65\,000$ fewer bacteria for the \omterm{def:b3:innate-immunity:innate}{neutrophils} to face.
Cost: $3\times 0.12\times 8 = \qty{2.9}{MJ}$ a day, a large meal. Worth
it when the pathogen is slowed and the host has reserves; not when
the pathogen grows as well at \qty{40}{\celsius}, or when the host is
starving, or when the fever itself approaches dangerous temperatures.
\end{solution}

\begin{solution}{exo:b3:innate-immunity:11}
Binding factor H recruits the host's own regulator: $d$ rises above
$a$ and the loop decays on the microbe as on a host cell. Cleaving
C3b removes deposited convertase: again a rise in $d$. A capsule
presents a surface on which convertase forms poorly, lowering $a$, and
hides the C3b that is deposited from phagocyte receptors. Cheapest:
one surface protein that binds factor H, which the host supplies free
and which does the work.
\end{solution}

\begin{solution}{exo:b3:innate-immunity:12}
At a splinter the dilation, leak and recruitment involve a few
millilitres of vessels and deliver defence where it is needed; the same
mediators released into the whole circulation dilate every vessel and
leak every capillary, so the pressure collapses and clotting is
triggered everywhere --- the physiology is identical, the scale is
lethal. TNF drives the synovial \omterm{def:b3:innate-immunity:inflammation}{inflammation} of rheumatoid arthritis,
so blocking it relieves the disease; but TNF is also what keeps
\omterm{def:b3:innate-immunity:innate}{macrophages} organised in the granulomas that contain dormant tubercle
bacilli, and blocking it lets latent tuberculosis escape.
\end{solution}

\begin{solution}{pb:b3:innate-immunity:1}
\textbf{1.} One hour is $1.5$ doublings: $10^{4}\times 2^{1.5} \approx
2.8\times 10^{4}$.
\textbf{2.} $2\times 10^{4}$ \omterm{def:b3:innate-immunity:innate}{neutrophils} kill up to $4\times 10^{5}$;
the bacteria grow only from $2.8\times 10^{4}$ to $8\times 10^{4}$:
killing outstrips growth fivefold.
\textbf{3.} Hour 2: $8\times 10^{4}$ before the kill, none after. The
population peaks at about $8\times 10^{4}$ near \qty{2}{h} and is zero
from then on; none remain at \qty{6}{h}.
\textbf{4.} Arrival at \qty{3}{h}: $10^{4}\times 2^{4.5} = 2.3\times
10^{5}$ then. Hour 4: $6.4\times 10^{5} - 4\times 10^{5} = 2.4\times
10^{5}$; hour 5: $6.8\times 10^{5} - 4\times 10^{5} = 2.8\times 10^{5}$;
hour 6: $7.9\times 10^{5} - 4\times 10^{5} = 3.9\times 10^{5}$ and
rising --- the \omterm{def:b3:innate-immunity:innate}{neutrophils} no longer keep up, and an abscess forms. Two
hours' delay turns a cure into a chronic infection.
\textbf{5.} $8\times 10^{4}/20 = 4000$ \omterm{def:b3:innate-immunity:innate}{neutrophils} die killing, plus the
$2\times 10^{4}$ that arrived and die anyway within a day. Pus is dead
\omterm{def:b3:innate-immunity:innate}{neutrophils}, dead bacteria and liquefied tissue; \omterm{def:b3:innate-immunity:innate}{macrophages} clear it
by eating the corpses, or it must drain.
\textbf{6.} Arrivals $2000$ an hour, kill $4\times 10^{4}$: hour 2,
$8\times 10^{4} - 4\times 10^{4} = 4\times 10^{4}$; hour 3, $1.1\times
10^{5} - 4\times 10^{4} = 7\times 10^{4}$; hour 4, $2\times 10^{5} -
4\times 10^{4} = 1.6\times 10^{5}$ --- growing without limit. The
patient cannot contain a trivial inoculum, so \omterm{def:b3:bacteriology:antibiotics}{antibiotics} must do the
killing from the first sign.
\textbf{7.} $B = 12.5(e^{0.08t} - 1)$: \qty{30}{s}, $125$; \qty{60}{s},
$1500$; \qty{120}{s}, $1.8\times 10^{5}$.
\textbf{8.} $e^{0.08t} = 1.6\times 10^{5}$: $t = 12/0.08 = \qty{150}{s}$.
\textbf{9.} $k_{0}/(d - a) = 10$.
\textbf{10.} $B(120) = 33(e^{3.6} - 1) \approx 1200$, $150$ times fewer.
The capsule buys minutes --- time to multiply, and protection until
antibodies arrive to start the classical pathway on the capsule
itself.
\textbf{11.} $10^{4}\times 2\times 10^{6} = 2\times 10^{10}$ C3b: one
millionth of the plasma's C3.
\textbf{12.} Without C3 all three pathways stop at their common step:
no \omterm{def:b3:innate-immunity:complement}{opsonisation}, no \omterm{def:b3:innate-immunity:complement}{anaphylatoxins}, no lysis --- severe recurrent
infections with encapsulated bacteria. Without C9 only the terminal
pore is lost; \omterm{def:b3:innate-immunity:complement}{opsonisation} and \omterm{def:b3:innate-immunity:inflammation}{inflammation} work, and the phenotype is
recurrent \emph{Neisseria} infections, the one genus that phagocytes
control poorly and lysis controls well.
\textbf{13.} $2\times 0.24\times 8 = \qty{3.8}{MJ}$, about \qty{100}{g} of
fat.
\textbf{14.} One doubling in the hour: $2\times 10^{4}$ instead of
$2.8\times 10^{4}$, a factor $1.4$.
\textbf{15.} Hour 2: $4\times 10^{4}$ before the kill, none after --- gone
at \qty{2}{h}, as before; fever's contribution matters when the
\omterm{def:b3:innate-immunity:innate}{neutrophil} supply is marginal, as in questions 4 and 6.
\textbf{16.} $199\,\text{mg/L}\times 3\,\text{L} = \qty{0.6}{g}$; $0.6/
\num{115000} = 5.2\times 10^{-6}$ mol $= 3.1\times 10^{18}$ molecules;
over \qty{172800}{s}, $1.8\times 10^{13}$ a second.
\textbf{17.} They lower the set point back toward \qty{37}{\celsius}, the
patient feels better, and the bacteria double a little faster; in this
model the \omterm{def:b3:innate-immunity:innate}{neutrophils} still win, and the evidence that antipyretics
worsen ordinary infections is weak.
\textbf{18.} Above about \qty{41}{\celsius} proteins begin to denature,
enzymes fail and neurons misfire (seizures); \qty{42}{\celsius} is
lethal. The set point is held below that by cryogens --- interleukin-10,
vasopressin, glucocorticoids --- and by the ceiling of prostaglandin
action in the hypothalamus.
\textbf{19.} $10^{9}\times 3\times 10^{6} = 3\times 10^{15}$ molecules
$= 5\times 10^{-9}$ mol in \qty{5}{L}: \qty{1e-9}{mol/L}.
\textbf{20.} Equal to the saturating concentration: every \omterm{def:b3:innate-immunity:innate}{macrophage} in
the body is maximally activated at once.
\textbf{21.} $10^{11}\times 1000\times 3600 = 3.6\times 10^{17}$ molecules
$= 6\times 10^{-7}$ mol; in \qty{15}{L}, \qty{4e-8}{mol/L} --- four
thousand times the active concentration.
\textbf{22.} Every arteriole dilates and every venule leaks: the blood
volume pools in the widened vessels and escapes into the tissues, so
the pressure falls; the kidneys, underperfused, stop filtering; and
clotting triggered in thousands of capillaries blocks them and consumes
the clotting factors.
\textbf{23.} Killed \omterm{def:b3:bacteriology:envelope}{Gram-negatives} release their whole \omterm{def:b3:bacteriology:envelope}{lipopolysaccharide}
at once, a bolus for TLR4, so the \omterm{def:b3:innate-immunity:inflammation}{cytokine} surge can worsen for hours.
Anti-TNF fails because by the time of shock the cascade has passed TNF
--- interleukin-1, interleukin-6 and coagulation are already running
--- and because TNF is needed to control the bacteria: timing, not
the target, is the problem.
\textbf{24.} Injected \omterm{def:b3:bacteriology:envelope}{lipopolysaccharide} does nothing to the TLR4-null
mouse, which survives doses lethal to normal mice; a live \omterm{def:b3:bacteriology:envelope}{Gram-negative}
infection kills it, because it cannot detect the bacteria early and
mount the \omterm{def:b3:innate-immunity:inflammation}{inflammation} that contains them. The paradox dissolves: the
detector that causes the shock is the detector that provides the
defence.
\textbf{25.} No bacteria remain at \qty{6}{h} (all killed by \qty{2}{h});
about $1.8\times 10^{5}$ C3b per bacterium at \qty{120}{s}; two days'
fever costs \qty{3.8}{MJ}, some \qty{100}{g} of fat.
\end{solution}
